\setcounter{section}{2}
\setcounter{theorem}{22}
\section{Quantum fine-grained complexity}

In this section, we give the formal definitions of the quantum fine-grained reduction and quantum strong exponential time hypothesis (QSETH). Moreover, we show that under QSETH, for $d=\polylog(n)$, the lower bounds for $\CP_{n,d}$ and $\OV_{n,d}$ are $n^{1-o(1)}$, which nearly matches the upper bounds given in \cref{cor:trivial_alg}.  

\subsection{Quantum fine-grained reduction and QSETH}

QSETH is defined based on the assumption that the best quantum algorithm for \textsf{CNF-SAT} is Grover search when the clause width $k$ is large enough. 
\begin{hypothesis}[QSETH]
For every $\epsilon > 0$, there exists a $k = k(\epsilon) \in \N$ such that no quantum algorithm can solve $k$-$\SAT$ (i.e., satisfiability on a CNF of width $k$) in $O(2^{(1/2 - \epsilon)n})$ time where $n$ is the number of variables. Moreover, this holds even when the number of clauses is at most $c(\epsilon) n$ where $c(\epsilon)$ denotes a constant that depends only on $\epsilon$.
\end{hypothesis}

Obviously, the Grover search can solve \textsf{CNF-SAT} in $\tilde{O}(2^{n/2})$. To the best of the our knowledge, there is no quantum algorithm that can do better than $O(2^{n/2})$ for any $k$. 


We recall that in the quantum query model, the input of a problem is given by a quantum oracle. Specifically, let $\mathsf{P}$ be a problem, and $X$ be an instance of $\mathsf{P}$ in the classical setting. Then, in the quantum query model, $X$ will be given by an oracle $\ora_X$. We will denote an algorithm or an oracle $\A$ with access to $\ora_X$ by $\A(\ora_X)$. 

We say $\A_{\epsilon}$ is an $\epsilon$-oracle for problem $\mathsf{P}$, if for every instance $\ora_X$, it holds that
\begin{align}
    \Pr[\A_{\epsilon}(\ora_X) = \mathsf{P}(X)] \geq 1-\epsilon,
\end{align}
and the running time is $O(1)$, where $\mathsf{P}(X)$ is the answer of $X$ for problem $\mathsf{P}$. 

\begin{definition}[Quantum oracles]\label{def:oisgood}
Let $X:=\{x_1,\dots,x_n\}$ be an instance of some problem and $\ora_X$ be the corresponding quantum oracle. To realize $\ora_X$, we do not need to write down the whole $X$; instead, we can just design a quantum circuit to realize the mapping 
\begin{align}
    \ket{i}\ket{0} \xrightarrow{\ora_X} \ket{i}\ket{x_i}.
\end{align}  
\end{definition}

\begin{definition}[Quantum fine-grained reduction]\label{def:qfgr}
Let $p(n)$ and $q(n)$ be nondecreasing functions of $n$. Let $\mathsf{P}$ and $\mathsf{Q}$ be two problems in the quantum query model and $\A_{\epsilon}$ be an $\epsilon$-oracle for $\mathsf{Q}$ with error probability $\epsilon\leq 1/3$. $\mathsf{P}$ is quantum ($p,q$)-reducible to $\mathsf{Q}$, denoted as $(\mathsf{P},p)\leq_{\QFG} (\mathsf{Q},q)$, if for every $\epsilon$, there exits a $\delta>0$, and algorithm $R$ with access to $\A_{\epsilon}$, a constant $d$, and an integer $k(n)$, such that for every $n\ge 1$, the algorithm $R$ takes any instance of $\mathsf{P}$ of size $n$ and satisfies the following: 
\begin{itemize}
    \item $R$ can solve $\mathsf{P}$ with success probability at least $2/3$ in time at most $d\cdot p(n)^{1-\delta}$.
    \item $R$ performs at most $k(n)$ quantum queries to $A_\epsilon$. Specifically, in the $j^{th}$ query, let $\mathbf{X}_j:= \{X_{1,j},X_{2,j},\dots\}$ be a set instances of $\mathsf{Q}$. Then, $R$ realizes the oracles $\{\ora_{X_{1,j}},\ora_{X_{2,j}},\dots\}$ in superposition and applies $A_{\eps}$ to solve the instances.
    \item The following inequality holds.
    \begin{align*}
        \sum_{j=1}^{k(n)}c(\mathbf{X}_j)\cdot q(n_j)^{1-\epsilon} \leq d\cdot p(n)^{1-\delta}, 
    \end{align*}
    where $c(\mathbf{X}_j)$ is the time required for $R$ to realize the oracles $\{\ora_{X_{1,j}},\ora_{X_{2,j}},\dots\}$ in superposition and $n_j := \max_{i}|X_{i,j}|$. 
\end{itemize}
\end{definition}

In \cref{def:qfgr}, the input of $A_{\eps}$ is given as a quantum oracle such that $A_{\eps}$ can be a quantum query algorithm with running time strictly less than the input size. Moreover, the quantum reduction $R$ can realize quantum oracles $\{\ora_{X_{1,j}},\ora_{X_{2,j}},\dots\}$ in superposition, and thus the time required is $\max_i c(X_{i,j})$ (where $c(X_{i,j})$ is the time required to realize $\ora_{X_{i,j}}$) instead of $\sum_{i} c(X_{i,j})$. This also allows $R$ to use fast quantum algorithms to process the information of $A_{\epsilon}'s$ output (e.g., amplitude amplification). 

\subsection{Lower bounds for $\CP$, $\OV$, and $\BCP$ in higher dimensions under QSETH}

Here, we give nearly linear lower bounds for $\OV$ and $\CP$ under QSETH by showing that there exist quantum fine-grained reductions from $\textsf{SAT}$ to these problems. 
\begin{theorem}\label{thm:qseth}
Assuming QSETH, for all $\epsilon>0$, there exists a $c$ such that $\OV_{n,c\log n}$ and $\CP_{n,(\log n)^c}$ cannot be solved by any quantum algorithm in time $O(n^{1-\epsilon})$.
\end{theorem}

We prove \cref{thm:qseth} by showing that there exist quantum fine-grained reductions from \textsf{CNF-SAT} to $\OV$, $\OV$ to $\BCP$, and $\BCP$ to $\CP$ with desired parameters.  We first give the reduction from \textsf{CNF-SAT} to $\OV$ as a warm-up. 

\begin{lemma}\label{lem:sat_q_ov}
\begin{align}
    (\mathsf{CNF\text{-}SAT}_{n},2^{n/2})\leq_{\QFG} (\OV_{n_1,d_1},n_1),
\end{align}
where $n_1=2^{n/2}$ and $d_1=\Theta(n)$.
\end{lemma}
\begin{proof}
Let $\phi$ be a CNF formula with $n$ variables and $m=\Theta(n)$ clauses. Let $\mathcal{A}$ be an algorithm for $\OV$. We first recall the classical reduction. Let $\phi := \phi_1\wedge \cdots \wedge\phi_m$. 
We divide the $n$ variables into two sets $A$ and $B$ with $|A| = |B| = \frac{n}{2}$. Let $A:=\{x_1,\dots,x_{n/2}\}$ and $B:= \{x_{n/2+1},\dots,x_{n}\}$. We let $S_A:=\{a_1,\dots,a_{2^{n/2}}\}$ be all assignments to $A$ and $S_B:=\{b_1,\dots,b_{2^{n/2}}\}$ be all assignments to $B$. We describe two mappings $f_A: S_A\rightarrow \{0,1\}^m$ and $f_B:S_B\rightarrow \{0,1\}^m$ as follows: 
\begin{align}
    &f_A(a_i) = [\phi_1(a_i),\dots,\phi_m(a_i)]^T, \text{ and}\\
    &f_B(b_i) = [\phi_1(b_i),\dots,\phi_m(b_i)]^T,
\end{align}
where $\phi_j(a_i) = 0$ if $a_i$ is a satisfied assignment for $\phi_j$, and  $\phi_j(a_i) = 1$ otherwise; we define $\phi_i(b_i)$ in the same way. Let $F_A:= \{f_A(a_i):\; i\in[2^{n/2}]\}$ and $F_B:= \{f_B(b_i):\; i\in[2^{n/2}]\}$. Then, it is obvious that if there exist $v\in F_A$ and $u\in F_B$ such that $\langle v, u \rangle=0$, then $\phi$ is satisfiable. However, at first glance, this reduction with $O(2^{n/2})$ running time is not fine-grained since we require the cost of the reduction to be at most $2^{n(1-\delta)/2}$ for some $\delta>0$ by \cref{def:qfgr}, but writing down elements in $F_A$ and $F_B$ already takes $\Omega(2^{n/2})$. %

Nevertheless, as in \cref{def:oisgood}, a quantum fine-grained reduction only needs to realize the functions $f_A$ and $f_B$, which takes $O(mkn)$ time where $k$ is the width of clauses. This is much less than $O(2^{n(1-\delta)/2})$. More specifically,  $f_A$ and $f_B$ are oracles for $F_A$ and $F_B$, and for any quantum query to elements in $F_A$ or $F_B$, the reduction can implement oracles $f_A$ and $f_B$:   
\begin{align}
    \ket{e,x}\ket{0}\xrightarrow{f_e} \ket{e,x}\ket{f_e(x)},
\end{align}
where $e\in \{A,B\}$, and the time $c(f_e)$ for the reduction to implement $f_e$ for one quantum query is at most $O(kmn)$. Finally, this reduction only uses one oracle ($F_A,F_B$). If there is an algorithm for $\OV$ which succeeds with probability $2/3$, we can boost the success probability of the reduction by repetition.  Therefore, ($\textsf{CNF-SAT},2^{n/2}$) is quantum reducible to ($\textsf{OV}_{n_1,d_1},n_1$). 
\end{proof}

Then, to prove $(\textsf{CNF-SAT}, 2^{n/2})\leq_{\QFG} (\textsf{CP}_{n_3,d_3}, n_3)$, we show that $(\textsf{BCP}_{n_2,d_2},n_2)\leq_{\QFG} (\CP_{n_3,d_3},n_3)$ and $(\textsf{OV}_{n_1,d_1}, n_1) \leq_{\QFG} (\textsf{BCP}_{n_2,d_2},n_2)$, where $n_2, n_3, d_2,d_3$ are some functions of $n$ specified in the following lemmas.

\begin{lemma}\label{lem:qbcp_to_cp}
For $d=\Theta(\log n)$,
\begin{align}
    (\BCP_{n,d},n)\leq_{\QFG} (\CP_{n',d'},n'),
\end{align}
where $n' = n^{O(1)}$ and $d'=(\log n)^{c}$ for some constant $c$ and all points have $\{0,1\}$ entries with the Hamming metric.
\end{lemma}

\begin{remark}
The points have coordinate entries in $\{0,1\}$, and the Hamming metric is equivalent to distance in $\ell_2$-metric (up to power of 2) in this case. Therefore, in the proof of Lemma~\ref{lem:qbcp_to_cp}, we can consider the Hamming distance between points instead of $\ell_2$ distance without loss of generality. 
\end{remark}

We first introduce the classical reductions in~\cite{r29} and some results we will use to prove Lemma~\ref{lem:qbcp_to_cp}. 

\paragraph{Classical reduction} 
We can consider an instance of $\BCP$ with two sets of points $A$ and $B$ as a weighted complete bipartite graph $K_{n,n}$, where the vertices are the points in these two sets and edges' weights are equal to the distances between the corresponding points. Then, solving $\BCP$ is equivalent to find an edge with the minimum weight in this graph.  
However, we cannot directly apply the algorithm for $\CP$ on this graph since there could be two points in the same set (no edge connecting them) that have a smaller distance than any pairs of points in two sets (connected by an edge). To overcome this difficulty, we can ``stretch'' the points to make the points in the same set far from each other, which is characterized by the contact dimension of a graph:
\begin{definition}[Contact Dimension]\label{def:contact_dim}
For any graph $G=(V,E)$, a mapping $\tau:V\rightarrow \R^d$ is said to realize $G$ if for some $\beta>0$, the following holds for every distinct vertices $u,v$:
\begin{align}
    \|\tau(u)-\tau(v)\|_2  =&~\beta ~ \text{if~} \{u,v\}\in E,\\
    \|\tau(u)-\tau(v)\|_2  > &~\beta ~ \text{otherwise.}\notag
\end{align}
The contact dimension of $G$, denoted by $\mathrm{cd}(G)$, is the minimum $d\in \N$ such that there exists $\tau:V\rightarrow \R^d$ realizing $G$.
\end{definition}
That is, with the help of $\tau$, we can restrict the optimal solution of $\CP$ to be the points connected by an edge in $G$. But we cannot realize the whole complete bipartite graph since $\mathrm{cd}(K_{n,n})=\Theta(n)$, which makes the dimension of the $\CP$ instance too large. \cite{r29} showed that we can realize a subgraph of $K_{n,n}$ and apply permutations to its vertices such that the union of these subgraphs cover $K_{n,n}$. In this way, $\BCP$ can be computed by solving $\CP$ on each subgraph and outputting the best solution. More specifically, the reduction in \cite{r29} relies on the following theorem:

\begin{theorem}[Theorem 4.2 in \cite{r29}] \label{thm:dense-cd}
For every $0<\delta < 1$, there exists a log-dense sequence $(n_i)_{i\in \N}$ such that, for every $i\in \N$, there is a bipartite graph $G_i=(A_i \dot \cup B_i, E_i)$ where $|A_i|=|B_i|=n_i$ and $|E_i|\ge \Omega(n_i^{2-\delta})$, such that $\mathrm{cd}(G_i)=(\log n_i)^{O(1/\delta)}$. Moreover, for all $i\in \N$, a realization $\tau:A_i\dot \cup B_i\rightarrow \{0,1\}^{(\log n_i)^{O(1/\delta)}}$ of $G_i$ can be constructed in $n_i^{2+o(1)}$ time.
\end{theorem}
\noindent The log-dense sequence is defined as follows:
\begin{definition}
A sequence $(n_i)_{i\in \N}$ of increasing positive integers is log-dense if there exists a constant $c\ge 1$ such that $\log n_{i+1}\le c\cdot \log n_i$ for all $i\in \N$.
\end{definition}
They also showed that, the permutations for covering the complete bipartite graph can be efficiently found, as shown in the following lemma.
\begin{lemma}[Lemma 3.11 in \cite{r29}] \label{lem:cover}
For any bipartite graph $G(A\dot \cup B, E_G)$ where $|A|=|B|=n$ and $E_G\ne \emptyset$, there exist side-preserving permutations $\pi_1,\dots, \pi_k:A\cup B\rightarrow A\cup B$ where $k\leq \frac{2n^2\ln n}{|E_G|}+1$ such that
\begin{align}
    \underset{i\in [k]}{\bigcup} E_{G_{\pi_i}} = E_{K_{n,n}}.
\end{align}
Moreover, such permutations can be found in $O(n^6\log n)$ time.
\end{lemma}

Now, we are ready to state the quantum fine-grained reduction by ``quantizing'' the classical reduction.
\begin{proof}[Proof of \cref{lem:qbcp_to_cp}]
Let $A,B$ be the two sets of input points of $\textsf{BCP}$. Suppose for \textsf{BCP}, there is an input oracle $\ora_{\BCP}$ which, given an index, returns the corresponding point:
\begin{align}
    \ket{b}\ket{i}\ket{0} \xrightarrow{\ora_{\BCP}}
    \begin{cases}
        \ket{b}\ket{i}\ket{x_i} &\text{if } b= 0,\\
        \ket{b}\ket{i}\ket{y_i} & \text{if } b= 1,
    \end{cases}
\end{align}
where $x_i$ is the $i$-th point in the set $A$ and $y_i$ is the $i$-th point in the set $B$. The sizes of $A$ and $B$ are both equal to $n$ and each point is in $\{0,1\}^{d_1}$, where $d_1=\Theta(\log n)$ is the dimension of $\BCP$. 
 
For $\CP$, suppose there is a quantum algorithm $\mathcal{A}$ such that for $m$ points in $\{0,1\}^{d_2}$ given by an oracle $\mathcal{M}_{CP}$, $\mathcal{A}^{\mathcal{M}_{CP}}$ returns the closest pair of these $n$ points with probability at least $2/3$.

Then we need to transform $\ora_{\BCP}$ to some oracles $\mathcal{M}_i$ for \CP, such that by running $\mathcal{A}$ with $\mathcal{M}_i$ as input oracles, we can get the bichromatic closest pair between $A$ and $B$. The reduction has four steps:

\paragraph*{1. Pre-processing.} We first follow the classical reduction to pre-process the input points of \textsf{BCP}. For some integer $n'\leq n^{0.1}$, we can partition $A$ and $B$ into $n'$-size subsets:
\begin{align}
    A= &~ A_1~\dot\cup ~\cdots ~\dot \cup~ A_{r},\\
    B= &~ B_1 ~\dot\cup ~\cdots ~\dot \cup~ B_{r},\notag
\end{align}
where $r=\lfloor n/n'\rfloor$. Here, we assume that $n$ is divisible by $n'$. It follows that
\begin{align}\label{eq:bcp_bcp}
    \textsf{BCP}(A,B)=\min_{i,j\in [r]} \textsf{BCP}(A_i,B_j).
\end{align}

Then, we use the algorithm in \cite{r29} to construct $k$ mappings $f_1,\dots, f_k:[2n']\rightarrow \{0,1\}^{d'}$ such that 
\begin{align}\label{eq:bcp_cp}
    \textsf{BCP}(A_i,B_j)=\min_{t\in [k]} \textsf{CP}(f_t(A_i)\cup f_t(B_j)) ~~~ \forall i,j\in [\lfloor n/n'\rfloor].
\end{align}

More specifically, we pick $n'$ to be the largest number in a log-dense sequence that is smaller than $n^{0.1}$. Then, we apply~\cref{thm:dense-cd} to classically construct a bipartite graph $G(A\cup B,E)$ with $n'$ vertices in each side and a realization $\tau$.
By choosing $\delta = \epsilon/2$ in \cref{thm:dense-cd}, the graph $G$ has $|E|=\Omega(n'^{2-\epsilon/2})$ edges. And we can get $2n'$ 0/1-strings of length $(\log n')^{O(2/\epsilon)}$:
\begin{align}
    \tau^A_i=\tau(u_i)~~~\forall u_i\in A, \quad\text{ and }\quad
    \tau^B_i=\tau(v_i)~~~\forall v_i\in B.
\end{align}

In order to cover the complete bipartite graph, we run the classical algorithm (\cref{lem:cover}) to find $k$ permutations $\pi_1,\dots, \pi_k:[n']\rightarrow [n']$, where $k$ is a parameter to be chosen later.

Then, we can define the mappings as follows:  
\begin{align}\label{eq:fcn_cp}
    f_t(u)=
    \begin{cases}
        x_v\circ \left(\tau^A_{\pi_t(w)}\right)^{d+1} & \text{if }1\leq u\leq n'\\
        y_v\circ \left(\tau^B_{\pi_t(w)}\right)^{d+1} & \text{if }n'<u\leq 2n'
    \end{cases} ~~~\forall t\in [k], u\in [2n'],
\end{align}
where $\circ$ means string concatenation and $(s)^{d+1}$ denotes $d+1$ copies of the string $s$. For a point $p\in A_i\cup B_j$, $u\in [2n']$ is the index in this union-set,  $v\in [n]$ is the index in the ground set $A$ or $B$, and $w\in [n']$ is the index in the subset $A_i$ or $B_j$. Further, if $1\leq u\leq n'$, then $w:=u$; otherwise, $w:=u-n'$.

\paragraph{2. Oracle construction.}
For $i,j\in [r],t\in [k]$, we then construct the input oracle $\mathcal{M}_{i,j,t}$ for the problem $\textsf{CP}(f_t(A_i)\cup f_t(B_j))$. For a query index $u\in [2n']$, 
\begin{align}\label{eq:oracle_cp}
    M_{i,j,t}\ket{u}\ket{0}=\ket{u}\ket{f_t(u)}.
\end{align}

With the help of the input oracle $\ora_{\BCP}$, we can implement $\mathcal{M}_{i,j,t}$ in the following way: 
\begin{enumerate}
    \item Prepare an ancilla qubit $\ket{b}$ such that $b=1$ if $u>n'$. 
    \item Transform $\ket{u}$ to $\ket{v}$, the index of the point in $A$ or $B$, based on the value of $b$. Note that the index is unique. Hence, this transformation is unitary and can be easily achieved by a small quantum circuit.
    \item Query $\ora_{\BCP}$ with input $\ket{b}\ket{v}$. Assume $b=0$. Then, 
    \begin{align}
        \ket{b}\ket{v}\ket{0}\xmapsto{\ora_{\BCP}}\ket{b}\ket{v}\ket{x_v}.
    \end{align}
    \item Similar to the second step, the index $w$ of the point in $A_i$ and $B_j$ can be computed from $v$ by a unitary transformation:
    \begin{align}
        \ket{b}\ket{v}\ket{x_v}\mapsto \ket{b}\ket{w}\ket{x_v}
    \end{align}
    \item Since each $w$ corresponds to a unique string $\tau^A_{\pi_t(w)}$, we can attach $d+1$ copies of this string to the remaining quantum registers:
    \begin{align}
        \ket{b}\ket{w}\ket{x_v}\mapsto \ket{b}\ket{w}\ket{x_v}\ket{\left( \tau^A_{\pi_t(w)} \right)^{d+1}}.
    \end{align}
    \item By recovering $u$ from $w$, we get the final state:
    \begin{align}
        \ket{u}\ket{f_t(u)}=\ket{u}\ket{x_v,\left( \tau^A_{\pi_t(w)} \right)^{d+1}}.
    \end{align}
\end{enumerate}

\paragraph*{3. Query process}
By \cref{eq:bcp_bcp,eq:bcp_cp}, we have
\begin{align}\label{eq:bcp_all}
    \textsf{BCP}(A,B)=\min_{i,j\in [r],t\in [k]}\textsf{CP}(f_t(A_i)\cup f_t(B_j)).
\end{align}
Hence, we can use quantum minimum-finding algorithm in \cref{cor:trivial_alg} over the sub-problems to find the minimum solution. For each sub-problem, we can run the algorithm for \textsf{CP} with $\mathcal{M}_{i,j,t}$ as the input oracle.

\paragraph*{4. Post-processing.} In case that $n$ is not divisible by $n'$, let the remaining points in $A$ and $B$ be $A_{res}$, $B_{res}$, respectively. Then, we can use Grover search to find the closest pair between $A_{res}$ and $B$, and between $B_{res}$ and $A$. Then, compare the answer to the previously computed result and pick the smaller one.

\paragraph*{Correctness.} In this reduction, we do not change the constructions of the mappings $\{f_i\}_{i\in [k]}$. By~\cite{r29}, \cref{eq:bcp_all} is correct in the classical setting. Hence, it also holds in the quantum setting, and we can use Grover search to find the minimum solution. However, since the algorithm $\mathcal{A}$ for $\CP$  has success probability $2/3$, for each tuple $(i,j,t)\in [r]\times [r]\times [k]$, we need to run $\mathcal{A}^{\mathcal{M}_{i,j,t}}$ $O(\log n)$ times to boost the success probability to at least $1-\frac{1}{n}$. Then, by the union bound, the probability that all queries in the Grover search are correct is at least $99/100$. Hence, by \cref{thm:mim}, the overall success probability is at least $2/3$.  


\paragraph*{Running Time of the Reduction.} 

The running time of the pre-processing step consists of two parts: (1) constructing the graph $G$ and its realization $\tau$; (2) finding $k$ permutations. For the first part, by~\cref{thm:dense-cd}, it can be done in $n'^{2+o(1)}$ time. For the second part, we pick $k=O(\frac{2n'^2\log n'}{n'^{2-\epsilon/2}})=O(n'^{\epsilon/2}\log n')$, and by~\cref{lem:cover}, it can be done in $O(n'^6 \log n')$ time. Hence, the total running time of pre-processing step is $n'^{2+o(1)}+O(n'^6\log n')=\wt{O}(n^{0.6})$.

The oracle construction can be done ``on-the-fly''. More specifically, given the strings $\{\tau^A_i,\tau^B_i\}_{i\in [n']}$, and permutations $\{\pi_i\}_{i\in [k]}$, for each query index $u$, we can simulate the oracle $\mathcal{M}_{i,j,t}$ defined in \cref{eq:oracle_cp} in $c(\mathcal{M}_{i,j,t})=O(d_2)=(\log n')^{\Omega(1)}=\wt{O}(1)$ time.

In the query process, for each $\CP$ instance indexed by $(i,j,t)$, suppose $\mathcal{A}^{\mathcal{M}_{i,j,t}}$ gets the answer in time $q(n')=n'$. Moreover, for each time $\mathcal{A}$ querying the input oracle $\mathcal{M}_{i,j,t}$, we need to spend $c(\mathcal{M}_{i,j,t})$ time to simulate the oracle. And we also have $O(\log n)$ runs for each instance. Hence, the total running time for each $\CP$ is at most  
\begin{align}
    n'^{1-\epsilon}\cdot \wt{O}(1)\cdot O(\log n)=\wt{O}(n'^{1-\epsilon}).
\end{align}
Then, we use Grover's search algorithm over $r^2\cdot k$ instances, which can be done by querying $\wt{O}(\sqrt{r^2\cdot k})$ instances by \cref{thm:mim}. Therefore, for any $\epsilon > 0$, we have

\begin{align}\label{eq:bcp_qfgr}
    \wt{O}(\sqrt{r^2 k})~\cdot &~ q(n')^{1-\epsilon}\cdot c(\mathcal{M}_{i,j,t})\cdot O(\log n) =  ~ 
    \wt{O}(\sqrt{(n/n')^2 k}\cdot(n')^{1-\eps}) \\
    \leq &~ \wt{O}(n\cdot (n')^{-\eps})
    \leq \wt{O}(n\cdot n^{-\eps/2})
    \leq n^{1-\delta},
\end{align}
where the first inequality follows from $k=O(n'^{\epsilon / 2}\log n')$ as shown in \cite{r29} and the last inequality follows by setting $\delta = \epsilon / 10$.

For the post-processing step, the sizes of $A_{res}$ and $B_{res}$ are at most $n'$. The running time is 
\begin{align}
    O(\sqrt{n\cdot n'}\cdot \log n) \leq \wt{O}(n^{0.55}).
\end{align}

Therefore, for any $\epsilon>0$, there exists a $\delta>0$ such that the \cref{eq:bcp_qfgr} holds and the total reduction time is $O(n^{1-\delta})$. By \cref{def:qfgr}, $\BCP_{n,d_1}$ can be quantum fine-grained reduced to $\CP_{n,d_2}$. This completes the proof of this lemma. 
\end{proof}

Finally, we show that $(\textsf{OV}_{n,d},n)\leq_{\QFG} (\textsf{BCP}_{n,d'},n)$ by quantizing the reduction in~\cite{r29} following the same idea. 
\begin{lemma}\label{lem:ov_q_bcp}
For $d=\Theta(\log n)$, 
\begin{align} 
    (\OV_{n,d},n)\leq_{\QFG} (\BCP_{n,d'},n),
\end{align}
where $d'=\Theta(\log n)$.
\end{lemma}
\begin{proof}
For an $\OV$ instance with sets of vectors $A$ and $B$, let $\ora_{\OV}$ be the input oracle such that:
\begin{align}
    \ora_{\OV}\ket{i}\ket{0}=\begin{cases}
        \ket{i}\ket{a_i} & \text{if }i\in A,\\
        \ket{i}\ket{b_i} & \text{if }i\in B.
    \end{cases}
\end{align}
where $a_i,b_i\in \{0,1\}^d$.

Then, similar to the classical reduction, we can construct mappings $f_A, f_B:\{0,1\}^d\rightarrow \{0,1\}^{5d}$ such that
\begin{align}
    f_A(a_i)_{5j-4:5j}=\begin{cases}
        11000 & \text{if }a_i(j)=0\\
        00110 & \text{if }a_i(j)=1
    \end{cases}~~~\forall j\in [d],
\end{align}
and
\begin{align}
    f_B(b_i)_{5j-4:5j}=\begin{cases}
        10100 & \text{if }b_i(j)=0,\\
        01001 & \text{if }b_i(j)=1.
    \end{cases} ~~~\forall j\in [d].
\end{align}

By the classical reduction, we have
\begin{align}
    \textsf{OV}(A,B)=1 \text{ if and only if }\textsf{BCP}(f_A(A),f_B(B))=2d
\end{align}
under Hamming distance. 

Also, note that we can simulate the input oracle $\ora_{\BCP}$ by first querying the oracle $\ora_{\OV}$ to get the vector, then applying the corresponding mapping $f_A$ or $f_B$, which can be done in $c(\ora_{\BCP})=O(d)$ time. Let the running time of the algorithm for $\textsf{BCP}$ be $q(n)=n$. Then for any $\epsilon>0$,
\begin{align}
    q(n)^{1-\epsilon}\cdot c(\ora_{\BCP})=n^{1-\epsilon}\cdot \Theta(\log n)\leq n^{1-\delta}
\end{align}
for some small $\delta >0$. Hence, by \cref{def:qfgr}, $(\textsf{OV}_{n,d},n)\leq_{\QFG} (\textsf{BCP}_{n,d'},n)$.
\end{proof}
 
\begin{proof}[Proof of \cref{thm:qseth}]

We can prove the theorem by contradiction following \cref{lem:sat_q_ov}, \cref{lem:ov_q_bcp}, and \cref{lem:qbcp_to_cp}. Specifically, suppose that there exists an $\epsilon>0$, for all $d=\Theta(\log n)$, there exists a quantum algorithm which can solve $\OV$ in time $O(n^{1-\epsilon})$. Then, we can obtain a quantum algorithm for \textsf{CNF-SAT}, which runs in time $O(2^{n/2(1-\epsilon)})$ by \cref{lem:sat_q_ov}. This contradicts QSETH. The proof for $\CP$ is the same. 
\end{proof}

\subsection{Quantum lower bound for $\BCP$ in nearly-constant dimensions under QSETH}
A byproduct of the previous subsection is a quantum lower bound for $\BCP$ in higher dimensions (i.e., $d = \polylog(n)$) under QSETH (\cref{lem:ov_q_bcp}). In this subsection, we show that this quantum lower bound for $\BCP$ even holds for nearly-constant dimensions (i.e., $d = c^{\log^*(n)}$). The main result of this subsection is the following theorem.
\begin{theorem}\label{thm:qseth_2_bcp}
Assuming QSETH, there is a constant $c$ such that $\BCP$ in $c^{\log^* (n)}$ dimensions requires $n^{1-o(1)}$ time for any quantum algorithm.
\end{theorem}

We will ``quantize'' the results by Chen~\cite{r17} to prove this theorem. More specifically, we first show a quantum fine-grained self-reduction of $\OV$ from $\log n$ dimensions with binary entries to $2^{\log ^* (n)}$ dimensions with integer entries ($\ZOV$). Then, we give a quantum fine-grained reduction from $\ZOV$ to $\BCP$ in nearly-constant dimensions.

\begin{definition}[Integral Orthogonal Vector, $\ZOV$]
Given two sets $A,B$ of $n$ vectors in $\mathbb \Z^d$, find a pair of vectors $a\in A$ and $b\in B$ such that
$\ip{a}{b}=0$, where the inner product is taken in $\mathbb{Z}$.
\end{definition}

We use $\ZOV_{n,d}$ to denote $\ZOV$  with $n$ vectors of $d$ dimension in each set. We then recap a theorem in \cite{r17}:

\begin{theorem}[{\cite[Theorem 4.1]{r17}}]\label{thm:zov_map}
Let $b,\ell$ be two sufficiently large integers. There is a classical reduction $\psi_{b,\ell}:\{0,1\}^{b\cdot\ell}\rightarrow \Z^\ell$ and a set $V_{b,\ell}\subseteq \Z$, such that for every $x,y\in \{0,1\}^{b\cdot \ell}$,
\begin{align}\label{eq:zov_iff}
    \ip{x}{y}=0~\Leftrightarrow~\ip{\psi_{b,\ell}(x)}{\psi_{b,\ell}(y)}\in V_{b,\ell}
\end{align}
and
\begin{align}
    0\leq \psi_{b,\ell}(x)_i< \ell^{6^{\log^* (b)}\cdot b}
\end{align}
for all possible $x$ and $i\in [\ell]$. Moreover, the computation of $\psi_{b,\ell}(x)$ takes $\poly(b\cdot \ell)$ time, and the set $V_{b,\ell}$ can be constructed in $O\left( \ell^{O(6^{\log^* (b)}\cdot b)}\cdot \poly(b\cdot \ell) \right)$ time.
\end{theorem}

Note that the size of $V_{b,\ell}$ is at most $\ell^{2\cdot 6^{\log^* (b)}\cdot b+1}$. 
The following lemma gives a quantum fine-grained reduction from $\OV$ to $\ZOV$:
\begin{lemma}\label{lem:ov_2_zov}
For $d=\Theta(\log n)$,
\begin{align}
    (\OV_{n,d}, n)\leq_{\QFG} (\ZOV_{n,d'}, n).
\end{align}
where $d'=2^{O(\log^* n_2)}$.
\end{lemma}
\begin{proof}
Consider an $\OV_{n,d}$ with $d=c\cdot \log n$, where $c$ is an arbitrary constant. We choose $\ell:=7^{\log^* n}$ and $b:=d/\ell$. Then, we can apply~\cref{thm:zov_map} to get the mapping function $\psi_{b,\ell}$ and the set $V_{b,\ell}$. For each $v\in V_{b,\ell}$, we'll construct an instance of $\ZOV_{n,\ell+1}$ as follows:
\begin{enumerate}
    \item Let $\ket{i}$ be the input query index of $\ZOV_{n,\ell+1}$.  
    \item Query $\OV_{n,d}$'s input oracle $\ora_{\OV}$ and get the vector $\ket{i,x}$.
    \item Compute the mapping $\psi_{b,\ell}$ and get $\ket{i,x}\ket{\psi_{b,\ell}(x)}$.
    \item If $x\in A$, then attach 1 to the end of the register: $\ket{i,x}\ket{\psi_{b,\ell}(x),1}$. If $x\in B$, then attach $-v$ to the end: $\ket{i,x}\ket{\psi_{b,\ell}(x),-v}$.
    \item Use $\ora_{\OV}$ to erase $x$ and return the final input state $\ket{i}\ket{\psi_{b,\ell}(x),1}$ or $\ket{i}\ket{\psi_{b,\ell}(x),-v}$.
\end{enumerate}
For each instance, we can use the quantum oracle for $\ZOV_{n,\ell+1}$ to check the orthogonality. $\OV_{n,d}$ is YES if and only if there exists a YES-instance of $\ZOV_{n,\ell+1}$.

\paragraph*{Correctness.} The correctness follows from~\cref{eq:zov_iff}: 
\begin{align}
\ip{x}{y}=0 \Leftrightarrow \ip{\psi_{b,\ell}(x)}{\psi_{b,\ell}(y)}=v\in V_{b,\ell} \Leftrightarrow \ip{[\psi_{b , \ell}(x), ~ 1]}{[\psi_{b,\ell}(y), ~ -v]}=0. 
\end{align}

\paragraph*{Reduction time.} Note that for $\ell=7^{\log^* n}$ and $b=d/\ell$, we have:
\begin{align}
    \log\left( \ell^{O(6^{\log^* (d)}\cdot b)} \right) = &~ \log \ell \cdot O\left( 6^{\log^* (d)}\cdot (d/\ell) \right)\\
    =&~ O\left( \log^* (n)\cdot 6^{\log^* n} \cdot c\log n / 7^{\log^* n} \right)\\
    =&~ o(\log n).
\end{align}
This implies that $|V_{b,\ell}|\leq \ell^{2\cdot 6^{\log^* (b)}\cdot b+1} \leq n^{o(1)}$.
Hence, the number of $\ZOV_{n,\ell+1}$ instances is $n^{o(1)}$ and the running time for compute $V_{b,\ell}$ is $n^{o(1)}$. And for each input query, the oracle for $\ZOV_{n,\ell+1}$ can be simulated in $c(\ora_{\ZOV})=\poly(d)=\poly(\log n)$ time. We can show that for every $\epsilon>0$, if $\ZOV_{n,\ell+1}$ can be decided in $n^{1-\epsilon}$ time, then
\begin{align}
    \sum_{v\in V_{b,\ell}} n^{1-\epsilon}\cdot c(\ora_{\ZOV})=n^{o(1)}\cdot n^{1-\epsilon}\cdot \poly(\log n) \leq n^{1-\delta}
\end{align}
for some $\delta > 0$, which satisfies the definition of quantum fine-grained reduction (\cref{def:qfgr}).

Therefore, $\OV_{n,O(\log n)}$ is quantum fine-grained reducible to $\ZOV_{n,2^{O(\log^* (n))}}$.
\end{proof}

Then, we give a quantum fine-grained reduction from $\ZOV$ to $\BCP$:
\begin{lemma}\label{lem:zov_2_bcp}
For $d=2^{O(\log^* n)}$,
\begin{align}
    (\ZOV_{n,d}, n)\leq_{\QFG} (\BCP_{n,d'}).
\end{align}
where $d'=d^2+2$.
\end{lemma}
\begin{proof}
We remark here that this proof closely follows that for Theorem 4.3 in \cite{r17}. We nonetheless give it here as some details are different.

For an $\ZOV_{n,d}$ instance with $(k\cdot \log n)$-bit entries, we construct a $\BCP$ instance as follows:
\begin{enumerate}
    \item For $x\in A$, construct a vector $x'\in \Z^{d^2}$ such that $x'_{i,j}=x_i\cdot x_j$. Here, we index a $d^2$-dimensional vector by $[d]\times [d]$. Similarly, for $y\in B$, construct a vector $y'\in \Z^{d^2}$ such that $y'_{i,j}=-y_i \cdot y_j$.
    \item Choose $W:=(d^2+1)\cdot n^{4k}$. For each $x'$, construct a vector $x''\in \R^{d^2+2}$ such that 
    \begin{align}
        x''=\left[x',~ \sqrt{W-\|x'\|_2^2},~ 0\right].
    \end{align}  For each $y'$, construct a vector $y''\in \R^{d^2+2}$ such that 
    \begin{align}
        y''=\left[y',~ 0, ~ \sqrt{W-\|y'\|_2^2}\right].
    \end{align}
\end{enumerate}
Then, we claim that the $\ZOV$ instance is YES if and only if the $\BCP$ instance has the minimum distance $\leq \sqrt{2W}$.

\paragraph*{Correctness.} First note that $\|x'\|_2^2\leq d^2\cdot (2^{k\log n})^4=d^2\cdot n^{4k}$. Hence, $W-\|x'\|_2^2 > 0$ and $W-\|y'\|_2^2>0$. For any $x''$ and $y''$ in the new constructed instance of $\BCP$, we have
\begin{align}
    \|x''-y''\|_2^2=&~ \|x''\|_2^2+\|y''\|_2^2-2\cdot \ip{x''}{y''}\\
    =&~ 2\cdot W-2\cdot \ip{x'}{y'}\\
    =&~ 2\cdot W - 2\cdot \sum_{(i,j)\in [d]\times [d]}x_i\cdot x_j\cdot (-y_j\cdot y_j)\\
    =&~ 2\cdot W + 2\cdot (\ip{x}{y})^2.
\end{align}
Hence,
\begin{align}
    \ip{x}{y}=0~\Leftrightarrow~\|x''-y''\|_2^2= 2W.
\end{align}

\paragraph*{Reduction time.} We can see from the above description that the input mapping function is simple and can be computed by a small quantum circuit in $O(d^2)=O(2^{O(\log^* (n))})$ time. Hence, we have $c(\ora_{\BCP})=O(2^{O(\log^* (n))})$. Also, by~\cref{def:qfgr}, it's easy to check that this is indeed a quantum fine-grained reduction from $\ZOV$ to $\BCP$.
\end{proof}

Now \cref{thm:qseth_2_bcp} follows immediately from \cref{lem:ov_2_zov} and \cref{lem:zov_2_bcp}:
\begin{proof}[Proof of \cref{thm:qseth_2_bcp}]
Let $\eps > 0$ be some constant. Suppose we can solve $\BCP_{n,c^{\log^* (n)}}$ in $n^{1-\epsilon}$ time for all constant $c>0$. Then, by \cref{lem:ov_2_zov} and \cref{lem:zov_2_bcp}, we can also solve $\OV_{n,c'\log n}$ in $n^{1-\delta}$ time for some $\delta > 0$ and any $c'>0$. However, this contradicts QSETH by~\cref{thm:qseth}. Therefore, assuming QSETH, there exists a constant $c$ such that $\BCP_{n,c^{\log^*( n)}}$ requires $n^{1-o(1)}$ time.
\end{proof}
